\chapter{Probability Theory}

\section{Expectation}

\subsubsection{Finite Case}
Let $X$ be a random variable with a finite number of outcomes $x_i$
for $i=1,\ldots,n$ with probability $p_i$.  The {\bf expectation} 
of $X$ is defined as
\be
\E(X) \defe \sum_{i=1}^n p_i \; x_i.
\ee
Because the sum of all probabilities is identically 1, 
\be
\sum_{i=1}^n p_i = 1,
\ee
the expected
value is considered the {\bf weighted average}, where each $p_i$ is the 
{\bf weight}.

\begin{Hremark}
The weighted average is also known as the {\bf first moment}, which is 
analogous to the {\bf center of mass} of a rigid body.
\end{Hremark}

\subsubsection{Equal Likelihood}
For the special case where each probability is equally likely, 
\be
p_1 = p_2 = \cdots = p_n = \frac{1}{n},
\ee
the expectation of $x$ is specialized to the form,
\be
\E(X) = \frac{1}{n} \sum_{i=1}^n \; x_i.
\ee
and is referred to as the {\bf simple average}.

\begin{Hremark}
The expectation of $X$, also known as the weighted average of $X$, should
not be confused with the simple average of $X$.  The latter is a special
case of the former.
\end{Hremark}

\section{Covariance}

For two jointly distributed real-valued random variables $x$ and $y$, with
{\bf expected values}\footnote{Also known as {\bf mean values}.} 
$\E(x)$ and $\E(y)$, respectively, the {\bf covariance} is defined as
\begin{align}
\cov(x,y) \; \defe \; & \; \E\left[ 
(x - \E(x)) \; (y - \E(y))
\right] \\
 & \; \E(xy) - \E(x) \E(y) - \E(x) \E(y) + \E(x) \E(y) \\
 & \; \E(xy) - \E(x) \E(y).
\end{align}

\subsubsection{Discrete Random Variables}
Let the random variable pair $(x,y)$ take on values $(x_i, y_i)$ for
$i=1,\ldots,n$ with probability $p_i$.  The covariance is defined as
\be
\cov(x,y) \defe \sum_{i=1}^n \; p_i \; \left(
x_i - \E(x)
\right)
\;
\left(
y_i - \E(y)
\right).
\label{eq:cov_discrete}
\ee
If the probability of each outcome is equal, the $p_i = 1/n$, and
(\ref{eq:cov_discrete}) becomes
\be
\cov(x,y) = \frac{1}{n} \; \sum_{i=1}^n \; \left(
x_i - \E(x)
\right)
\;
\left(
y_i - \E(y)
\right).
\label{eq:cov_discrete_equal}
\ee

\subsection{Pearson Correlation Coefficient}
The {\bf Pearson correlation coefficient} (PCC)\footnote{Also known as the
Pearson's $r$ coefficient, Pearson product-moment correlation
coefficient (PPMCC), and the bivariate correlation.} is a 
measure of the linear correlation between two variables 
$X$ and $Y$.  

When a pair of random variables $(X,Y)$ are applied to a population, the PCC
is called the {\bf population correlation coefficient} and denoted as 
$\rho_{X,Y}$ and is defined as 
\be
\rho_{X,Y} = \frac{\cov(X,Y)}{\sigma_X \; \sigma_Y}, 
\ee
where 
\begin{itemize}
\item \cov is the covariance of $X$ and $Y$, 
\item $\sigma_X$ is the standard deviation of $X$, and 
\item $\sigma_Y$ is the standard deviation of $Y$.
\end{itemize}
The PCC is bounded as
\be
-1 \leq \rho_{X,Y} \leq 1.
\ee
\begin{itemize}
\item A value of -1 indicates complete negative linear correlation.
\item A value of  0 indicates the absence of linear correlation.
\item A value of +1 indicates complete positive linear correlation.
\end{itemize}

\chapter{Modern Portfolio Theory}
% See https://en.wikipedia.org/wiki/Outline_of_finance#Portfolio_theory
% Also Bessel's correction for sample means and deviations.
% https://en.wikipedia.org/wiki/Bessel%27s_correction

\section{Cross-Correlation}
In signal processing, {\bf cross-correlation}\footnote{Cross-correlation is
also called the {\bf sliding dot product} or {\bf sliding inner product}}
is a measure of similarity of two signals.  

In probability theory, cross-correlation refers to the correlations between
two random variables $X$ and $Y$.


\section{Portfolio and Assets}

Let a portfolio $P$ be a collection of $n$ individual assets 
with present values $a_1, a_2, \ldots a_n$.
Assets have units of 
U.S.~dollars (USD).\footnote{When the asset is stock, bond, or 
other instrument that consists of 
fractional ownership (shares) denominated by price per share, 
then the present value of the asset is simply the number of shares
held (holdings) in the portfolio multiplied by the price per share of that
holding.}

The present valuation of the portfolio, $PV(P)$, is given as
\be
 PV(P) = \sum_{i=1}^n a_i.
\ee

\section{Portfolio Weighting}

The relative weight $w_i$ of each asset $a_i$ is simply the fraction 
of the asset's present value to the portfolio's present value, thus
\be
 w_i = \frac{a_i}{PV} = \frac{a_i}{\sum_{i=1}^n a_i}.
\ee
The sum of all weights, must equal one (1).
\be
 \sum_{i=1}^n w_i = \sum_{i=1}^n \frac{a_i}{PV} = \frac{1}{PV} \sum_{i=1}^n a_i 
 = \frac{PV}{PV} = 1.
\ee

\section{Portfolio Return}

Let $\E(R_P)$ and $\E(R_i)$ the be expected periodic return $R_P$ 
of the portfolio $P$ and 
expected periodic return $R_i$ on an asset $a_i$, 
respectively.\footnote{The period
is typically stated annually, quarterly, or yearly.}
The expected portfolio return is then simply the weighted sum of each 
asset's return, 
\be
\E(R_P) = \sum_{i=1}^n w_i \; \E(R_i).
\ee

Let the mean periodic return $\bar{R}(a_i)$ of asset $a_i$ is
given by observing the return at $K$ times during the periodic return 
interval, 
\be
\bar{R}(a_i) = \frac{1}{K} \sum_{k=1}^K R_k(a_i).
\ee
For clarity, the notation $f(a_i)$ to denote {\em on a per-asset basis} 
is assumed, as written hereinafter simply as $f$.
Then the mean periodic return is given as
\be
 \bar{R} = \frac{1}{K} \sum_{k=1}^K R_k.
\ee
Let $\var(R) = s^2$ be the (sample) variance of periodic returns on an 
asset $a_i$, given by
\be
 \var(R) = s^2 = \frac{1}{K-1} \sum_{k=1}^K \left(R_k - \bar{R}\right)^2.
\ee
The (sample) standard deviation $s$ is the positive square root of the 
(sample) variance $s^2$,
\be
 s = \sqrt{\frac{1}{K-1} \sum_{k=1}^K \left(R_k - \bar{R}\right)^2}.
\ee


